\documentclass[11pt]{article}
\usepackage[margin=2.2cm]{geometry}
\usepackage{enumitem}
\usepackage{amsmath}
\usepackage{amssymb}
\usepackage{titlesec}
\usepackage{hyperref}
\usepackage[numbers,sort&compress]{natbib}

\titlespacing*{\section}{0pt}{10pt}{4pt}
\titlespacing*{\subsection}{0pt}{8pt}{3pt}
\setlist{nosep, leftmargin=1.2em}

\title{State of the Art \\ \large LLM-Based Persuasive User Guidance for Socially Optimal Mobility Decisions}
\author{Fares Memmi \\ \small Supervisors: Angelo Furno, Rim Slama Salmi --- EMob-Lab, Universit\'e Gustave Eiffel / ENTPE}
\date{}

\begin{document}
\maketitle

\section*{Introduction}
This project studies whether a conversational, LLM-based advisor can persuade players in a repeated route-choice game toward the system-optimal traffic distribution, and whether a more structured advisor architecture persuades better than a single fixed-prompt baseline. This places the work at the intersection of three traditionally separate bodies of literature: (i) traffic assignment and route-choice theory, which defines what ``optimal'' means and how individual choices aggregate into collective outcomes; (ii) human--AI interaction and trust, which governs how and when users follow AI advice; and (iii) LLM-based agentic and persuasive systems, which supplies the mechanism by which advice is generated and delivered. This chapter reviews each of these areas in turn, together with serious games as the experimental methodology that ties them together, and closes by positioning the gap this project addresses.

\section{Traffic Assignment and Route Choice}
Classical traffic assignment theory formalizes how individual route choices aggregate into network-level travel times. Wardrop's first principle characterizes \emph{user equilibrium}: at equilibrium, no traveler can reduce their own travel time by unilaterally switching routes, so all used routes between an origin--destination pair have equal (and minimal) travel time \citep{wardrop1952}. This is distinct from the \emph{system optimum}, the flow allocation that minimizes total travel time across all users, which generally requires some users to accept longer routes for the collective benefit. The gap between user equilibrium and system optimum --- the ``price of anarchy'' of selfish routing --- is precisely the quantity this project's advisor is designed to close: the system-optimal split computed in the game (via brute-force search over route splits) is the same theoretical object as Wardrop's system optimum, and the ``gap-to-optimal'' metric used to evaluate PersuLLM-1 and RegConSuader is a direct operationalization of this classical gap.

Flow-dependent cost functions, most commonly of Bureau of Public Roads (BPR) form, are the standard way to map route flow to realized travel time in a tractable macroscopic model, and are foundational to equilibrium computation and traffic assignment more broadly \citep{sheffi1985}. The project's own congestion model --- $t = t_{\text{free}} \times (1 + 0.15 \cdot (\text{flow}/\text{capacity})^4)$ --- is a direct instance of this formulation, chosen for the same reasons the classical literature adopts it: analytical tractability, a convex and sharply increasing cost near capacity, and ease of controlled experimentation, at the cost of the fidelity offered by microscopic simulators such as SUMO or MATSim.

At the individual level, discrete choice theory explains how travelers select among a finite set of alternatives (here, three generated routes) based on perceived attributes such as travel time and reliability \citep{benakiva1999}. These models typically assume static, non-interactive decision processes: an agent observes attributes and applies a fixed choice rule. They do not, by construction, represent a traveler who discusses the decision with an advisor, revises beliefs through dialogue, or adapts a decision rule across repeated rounds based on trust built with an AI agent. This is precisely the first gap the project addresses: replacing (or augmenting) a static choice rule with a language-mediated, iterative interaction.

\section{Traffic Prediction and Decision-Support Systems}
Short-term traffic forecasting has been studied extensively, from classical statistical and autoregressive models to deep learning approaches capable of capturing nonlinear spatio-temporal patterns in large-scale traffic data \citep{vlahogianni2014, lv2015}. These methods focus on predictive accuracy of travel times or flows, typically as an end in itself or as an input to static route-guidance systems (e.g., navigation apps recommending the fastest route).

The common thread across this literature is that the interface between prediction and user remains a passive display: a number, a color-coded map, or a ranked list of routes. The user is not invited to question the prediction, ask about its uncertainty, or negotiate a trade-off with the system. This project's platform retains a lightweight predictive layer (predicted travel time as a function of prior realized times, flows and context) but replaces the passive display with a conversational advisor capable of explaining, contextualizing and justifying its recommendation --- turning a one-way forecast into a two-way interaction, in line with the motivation identified in the original internship proposal for moving beyond static recommender outputs \citep{vlahogianni2014}.

\section{Serious Games as an Experimental Methodology}
A serious game is a game designed with an explicit purpose beyond entertainment --- training, education, or, as here, controlled scientific observation \citep{abt1970, michaelchen2005}. The format is well suited to studying repeated decision-making because it allows a realistic but stylized decision situation to be replayed under controlled conditions, with measurable feedback loops and quantifiable strategic adaptation across rounds, while remaining tractable and reproducible in a way that field studies of real traffic behavior are not.

This project's platform follows this logic directly: a 5$\times$5 grid network with five distinct scenarios (clean baseline, spatial congestion, a bottleneck trap, network disruption, and an express-corridor risk/reward setup) provides a controlled but varied set of repeated decision situations, each round constituting a decision epoch rather than a physical time interval. The round structure (predicted time shown, choice submitted, realized time computed from aggregate choices, feedback given, next round begins from the updated state) mirrors the closed-loop repeated-game design described in the original proposal, and is the concrete instrument used to measure how persuasion changes behavior over the five scenarios rather than in a single one-shot decision.

\section{Human--AI Interaction and Trust in Automated Advice}
A central concern in human--AI interaction is reliance: under what conditions do users follow, ignore, or override automated recommendations, and how does this reliance evolve with repeated exposure to advice that is sometimes right and sometimes wrong? The broader ``machine behaviour'' perspective argues that AI systems' interactions with humans and with each other need to be studied empirically as a behavioral phenomenon in its own right, not simply inferred from the systems' design \citep{rahwan2019}. This motivates treating compliance, trust and interaction patterns in this project as empirical outcomes to be measured (via compliance rate, a trust moving-average over rounds, and interaction intensity) rather than assumed.

Recent work has begun to study this reliance dynamic specifically with LLM-based advisors in decision-support settings. Ma et al. examine \emph{deliberative} AI-assisted decision-making, where an LLM does not simply hand down a recommendation but engages the user in a structured back-and-forth intended to surface reasoning and calibrate trust \citep{ma2025deliberation}. Wang et al. directly compare AI-driven and human decision-making in a repeated daily-commute setting, finding that LLM agents can learn human-like commuting strategies but tend to underperform both humans and reinforcement-learning baselines, and exhibit distinctive failure modes such as GPT-3.5's excessive stability versus GPT-4's excessive volatility, alongside difficulty handling collaborative, multi-agent dynamics \citep{wang2025comparing}. This is directly relevant to the project's own risk, flagged explicitly in the research hypotheses (H2--H5), that AI-mediated coordination can improve average efficiency while introducing new failure modes such as herding or oscillation --- a risk this project measures with an entropy-based route-concentration metric and a round-to-round flow-stability metric.

\section{LLM-Based Agentic Systems}
The LLM agentic layer in this project's architecture --- an advisor that accesses structured state, reasons about it, and produces a contextualized recommendation --- follows the general pattern established by ReAct, which interleaves reasoning traces with actions (e.g., tool calls or information retrieval) so that an LLM's output is grounded in, and can be updated by, external state rather than produced from a static prompt alone \citep{yao2023react}. The project's own call structure implements a simplified version of this idea: each advisor call is stateless, but is grounded by a context block rebuilt from the database on every call (live routes, current distribution, the computed optimal split, and the player's own history), so that continuity of ``reasoning'' across rounds is achieved by re-injecting state as text rather than by persisting hidden state inside the model.

Two recent, more directly relevant works apply LLM agents specifically to route-choice and transportation recommendation. Wang et al. study \emph{agentic} LLMs for day-to-day route choice, modeling travelers as LLM-driven agents that adapt their route decisions over repeated days, which is architecturally close to this project's simulated population of persona-driven agents used for cheap, repeatable pre-human validation \citep{wang2024agentic}. Afridi et al. propose an LLM-based recommender system for smart transportation built around the idea of \emph{serendipity} --- surfacing recommendations that are useful precisely because they are not the obvious default choice --- which is a relevant counterpoint to this project's advisor, whose recommendation is instead anchored to a single, well-defined optimality criterion (the system-optimal split) rather than serendipitous discovery \citep{afridi2024serendipity}.

\section{Computational and Persuasive AI}
Where the previous section covers agents that inform or recommend, this project's advisor is explicitly designed to \emph{persuade} --- to move behavior toward a target (the system-optimal split) that may not align with a user's naive individual incentive. This connects the project to the literature on computational persuasion. Cialdini's classic taxonomy of influence principles (reciprocity, commitment/consistency, social proof, authority, liking, scarcity, and unity) provides the vocabulary most directly adapted into this project: RegConSuader's strategy set (authority, social proof, and consistency) is a direct subset of these principles \citep{cialdini2021}.

The MA2P framework operationalizes this vocabulary inside an LLM-based agent architecture, proposing a meta-cognitive structure for complex persuasion tasks and explicitly drawing on Cialdini's principles as its strategic repertoire \citep{zhang2024ma2p}. MA2P is the direct architectural precedent for RegConSuader: RegConSuader's decomposition into a Configurator (estimating how strongly and in what tone to push), a Strategy Selector (choosing a persuasion principle), a Persuader (generating the actual discussion), and a Decider (regulating the exchange turn-by-turn) mirrors MA2P's separation of persuasion into distinct cognitive sub-tasks rather than a single monolithic prompt, and its choice of exactly three strategies (out of Cialdini's larger set) reflects a deliberate initial scoping decision rather than a claim that only these three are relevant. MA2P's own framing draws in turn on LeCun's proposal for autonomous machine intelligence structured as discrete perceive--act--predict modules \citep{lecun2022}, which is the conceptual justification, inherited via MA2P, for why RegConSuader is built as coordinated components rather than one larger prompt.

\section{Multi-User and Multi-Agent Human--AI Studies}
Most human--AI interaction research, including much of the work reviewed above, studies a single user interacting with a single system. Traffic, however, is inherently a shared-resource problem: one user's route choice changes the travel time experienced by every other user on that route. This project's central and personal-agent conditions are designed specifically to probe this multi-user dimension --- whether a single advisor with global visibility (central agent) coordinates better, and induces different collective effects (e.g., over-coordination, herding), than individually customized advisors with no explicit cross-user coordination (personal agents).

Wang et al.'s comparison of AI and human decision-making in \emph{daily collaborative} experiments is one of the few works directly addressing this multi-agent, repeated setting, and its finding that reinforcement-learning agents align more closely with human collaborative benchmarks than LLM agents do is a useful caution: it suggests LLM-based coordination advice cannot be assumed to improve collective outcomes by default, and must be measured empirically \citep{wang2025comparing}. The ``machine behaviour'' perspective similarly argues for treating multi-agent, human-AI-mixed systems as objects of direct empirical study rather than systems whose collective properties can be derived analytically from individual-agent design \citep{rahwan2019}.

\section*{Summary of the Gap}
Across these five areas, a consistent gap emerges. Traffic assignment theory defines the system-optimal target but assumes static, non-interactive choice rules; traffic prediction research optimizes forecast accuracy but treats the user interface as a passive display; human--AI trust research studies reliance on advice but rarely in a repeated, shared-resource setting; LLM agentic architectures provide the technical means for grounded, stateless-but-context-rich advisors but have seen limited application to explicitly \emph{persuasive}, target-directed guidance in transportation; and persuasive-AI frameworks such as MA2P supply a structured strategy repertoire but have not been evaluated against a socially optimal traffic target in a multi-user serious game. This project's contribution is to combine these strands: a classical system-optimal target, delivered through a structured, strategy-aware LLM advisor, evaluated in a repeated multi-user serious game against both a no-AI baseline and a minimal single-prompt persuader (PersuLLM-1), with explicit attention to the collective failure modes (herding, oscillation, fairness loss) that the multi-agent and human--AI-trust literature suggests are the real risk of AI-mediated coordination.

\bibliographystyle{unsrtnat}
\bibliography{references}

\end{document}
